\documentclass{article}
\usepackage{graphicx} % Required for inserting images
\usepackage{biblatex}
\addbibresource{references.bib}
\usepackage{amstext}

\title{Honors Physics}
\author{Fernando Delgado}
\date{August 2026}

\begin{document}
\maketitle
\section{Introduction}

\subsection{MINFLUX concept}
Single-molecule fluorescence localization with minimum photon flux imaging
(MINFLUX) can achieve localization precisions in the small nanometer range or better under
suitable conditions,Potentially adverse conditions, such as a fixed fluorescence dipole or optical
aberrations, that could cause systematic localization errors, have received little attention up to
now. \cite{stallingaAbsorptionDipoleEffects2026}
\subsection{single particle analysis}
\subsection{motivation requirements}
\section{EOD}
\subsection{physical concept}

An electro-optic deflector (EOD) steers a laser beam by changing the refractive index of a
solid-state crystal under an applied electric field. Because the deflection is produced by the
field rather than by a moving surface, the device contains no mechanical parts, and this single
fact sets both its advantage and its limitation. Solid-state deflectors reach very high angular
velocities but span only small deflection angles, where mirror-based mechanical scanners cover
large angles comparatively slowly
\cite{romerElectroopticAcoustoopticLaser2014,heberleElectroopticAcoustoopticLaser2016}. The two
families are usefully compared through the number of resolvable spots and the rate at which those
spots can be addressed; on the latter measure, electro-optic and acousto-optic deflectors exceed
conventional mechanical scanners by up to an order of magnitude
\cite{heberleElectroopticAcoustoopticLaser2016}.

The underlying index change is the electro-optic effect. In the quadratic, or Kerr, regime the
index shift scales with the square of the applied field and arises from the third-order
nonlinearity of the medium, in contrast to the linear electro-optic effect more commonly exploited
in optical crystals \cite{qasymehQuadraticElectroOpticKerr2008}. Potassium tantalate niobate,
KTa$_{1-x}$Nb$_x$O$_3$ (KTN), is attractive for deflector work because its electro-optic response
is unusually large, and crystals of optical quality have been grown specifically to build
high-speed beam deflectors and variable-focus lenses \cite{yagiElectroopticKTNDevices2014}.

The way an applied voltage becomes a deflection angle in KTN is less obvious than a simple index
gradient would suggest. A rectangular crystal carrying uniform electrodes, with none of the prism
geometry used in conventional electro-optic deflectors, still produces a large deflection: 250~mrad
($14.3^\circ$) has been measured from a 0.5~mm thick KTN crystal of only 5.0~mm interaction length
under $\pm250$~V \cite{nakamuraSpacechargecontrolledElectroopticEffect2008}. The mechanism is
space-charge-controlled conduction. Electrons injected from an Ohmic contact accumulate inside the
crystal and establish a non-uniform internal field, which through the electro-optic effect becomes
a refractive-index gradient across the beam path, and it is that gradient which bends the beam
\cite{nakamuraSpacechargecontrolledElectroopticEffect2008}. Independent analysis of KTN scanners
reaches the same conclusion, identifying trapped space charges as the principal cause of the
deflection and of its dependence on operating conditions \cite{bechtoldBeamProfileDeformation2012}.

This mechanism is fast. Electro-optic scanning above $10^4$~rad/s at deflection angles beyond
$\pm100$~mrad has been demonstrated \cite{bechtoldBeamProfileDeformation2012}, and the scan angle of
a KTN deflector has been shown to remain unchanged up to a drive frequency of 400~kHz
\cite{miyazu400KHzBeam2010}.

% TODO: crystal specification for the EOD used in this work — material, dimensions, and drive
% voltage range. This is instrument-specific and must come from the device datasheet or direct
% measurement; it cannot be sourced from the literature.

\begin{itemize}
    \item [FIGURE: diagram] schematic of EOD crystal + applied field + resulting beam deflection
\end{itemize}

\subsection{characterization}

\subsubsection{intro why}

MINFLUX determines an emitter's position by probing it with an intensity minimum rather than by
collecting the full diffraction pattern of its fluorescence, and this is what allows it to reach
single-digit nanometer precision while detecting ten to a hundred times fewer photons than
centroid-based localization
\cite{balzarottiNanometerResolutionImaging2017,eilersMINFLUXMonitorsRapid2018}. The same property
is the method's vulnerability. Because the position estimate is referred to the zero of the
excitation pattern, the quality of that zero is not a cosmetic concern but enters the measurement
directly, which is why the generation and optimization of toroidal foci is treated as a discipline
of its own \cite{lopezOptimizationCharacterizationToroidal2023}. Conditions that distort the
excitation pattern produce not merely a noisier estimate but a biased one: simulations of fixed
absorption dipoles and optical aberrations in MINFLUX find systematic localization errors reaching
roughly a quarter of the diameter of the circle spanned by the doughnut centre positions
\cite{stallingaAbsorptionDipoleEffects2026}. An error of that character cannot be averaged away by
collecting more photons, which is the reason the beam delivered by the EOD must be characterized
rather than assumed.

Speed matters for a related reason. The photon economy of MINFLUX is what opens the microsecond
regime, where molecular jumps and conformational changes become observable
\cite{eilersMINFLUXMonitorsRapid2018}, and realizing that regime on a conventional microscope stand
has depended specifically on synchronized electro-optical and galvanometric beam steering together
with sub-nanometer sample stabilization \cite{schmidtMINFLUXNanometerscale3D2021}. The EOD is
therefore not an incidental convenience in the beam path but one of the components that the
achievable time resolution rests on, and its response must hold at the rates the experiment demands.

Three failure modes bear on the measurements that follow. Beam degradation is documented for
electro-optic scanning in KTN, where the profile deformation depends on the midpoint of the
deflection angle, on the amplitude of the deflection, and on the laser fluence
\cite{bechtoldBeamProfileDeformation2012}; a deflection-dependent distortion of this kind would
make the excitation minimum a function of scan position, and so couple directly into the
localization. Deflection nonlinearity is the artifact that motivated the substitution of EODs for
acousto-optic deflectors in high-resolution force-clamp trapping, where the nonlinear response of
the AODs introduced experimental artifacts that switching to EODs removed
\cite{woodyElectroopticDeflectorsDeliver2018}; the concern transfers to MINFLUX, since a nonlinear
mapping from drive voltage to angle misplaces the excitation pattern relative to where the control
system believes it to be. Pointing drift is the third, and techniques that hold the sample fixed
within an excitation pattern are known to require active drift correction to approach their
theoretical limits at all \cite{lopezOpensourceSubNanometerStabilization2025}.

\subsubsection{setup}
\begin{itemize}
    \item beam degradation through alignment
    \begin{itemize}
        \item [EXPERIMENT] camera + EOD on optical rail, frames captured at multiple fixed positions along beam path (with/without EOD, before/after subsequent optics)
        \item [GRAPH] ideal Gaussian beam profile vs. measured beam after EOD (2D heatmap + fit overlay)
        \item [GRAPH] beam quality metrics vs. position along path (drift through alignment)
        \item [GRAPH] cross-correlation similarity (peak NCC) and beam displacement ($\Delta y$, $\Delta x$) vs. position
    \end{itemize}
    \item beam quality
    \begin{itemize}
        \item [EXPERIMENT] control (no-EOD) frame vs. EOD-deflected frame, same optical path
        \item [GRAPH] FWHM$_x$, FWHM$_y$ vs. position/condition
        \item [GRAPH] ellipticity vs. position/condition
        \item [GRAPH] $R^2$ (fit quality) vs. position/condition
        \item [GRAPH] beam radius $\omega$ (fit-derived vs. FWHM-derived) vs. position/condition
    \end{itemize}
    \item range angle
    \begin{itemize}
        \item [EXPERIMENT] HVA200 voltage sweep (DAQ output, stepped range, settle time), camera records beam centroid at each step
        \item [GRAPH] voltage vs. angle (calibration curve, mrad/V)
        \item [GRAPH] voltage vs. beam displacement (px / µm)
        \item [TABLE/TEXT] measured mrad/kV vs. manufacturer spec (% of spec)
    \end{itemize}
    \item polarization based angles
    \begin{itemize}
        \item [EXPERIMENT] polarization angle set/measured relative to EOD deflection axis
        \item [GRAPH] polarization angle vs. deflection-axis position (alignment check)
    \end{itemize}
\end{itemize}

\subsubsection{methods}
\begin{itemize}
    \item background subtraction (darkest 5–10\% pixels, subtracted before fitting)
    \item control-frame referencing (baseline / no-deflection reference for all comparisons)
    \item centroid localization: peak-pixel (live) vs. optimized 2D Gaussian fit (offline)
    \item frame averaging / repeat counts per condition

    \item beam degradation through alignment
    \begin{itemize}
        \item multi-position acquisition along optical path
        \item 2D Gaussian fit per position (Nelder--Mead least-squares)
        \item control $\omega$ used as shape reference for non-control positions
        \item 2D normalized cross-correlation vs. control (similarity, sub-pixel displacement)
    \end{itemize}

    \item beam quality metrics
    \begin{itemize}
        \item FWHM: half-max crossing of background-subtracted profile
        \item ellipticity: max/min FWHM ratio
        \item $R^2$: fit vs. measured frame goodness-of-fit
        \item beam radius $\omega$: from fit and independently from FWHM (cross-check)
        \item pixel $\to$ physical unit conversion (µm/px)
        \item fit initialization strategy (peak center, background amplitude, profile-based $\omega$)
    \end{itemize}

    \item range angle
    \begin{itemize}
        \item voltage sweep procedure (start, stop, steps, settle time)
        \item DAQ output + monitor readback recorded per step
        \item centroid displacement (px) $\to$ physical displacement (µm)
        \item small-angle approximation: $\theta = \arctan(\Delta x / L)$
        \item calibration: mrad/V, validated against manufacturer spec
    \end{itemize}

    \item polarization-based angles
    \begin{itemize}
        \item polarization angle set/measured independently
        \item correlation vs. deflection axis orientation
        \item alignment check: electro-optic axis vs. polarization axis
    \end{itemize}
\end{itemize}

\subsubsection{results}
\begin{itemize}
    \item beam degradation through alignment
    \begin{itemize}
        \item [GRAPH] measured vs. ideal Gaussian at each position (2D + profile overlay)
        \item [GRAPH] FWHM/ellipticity/$R^2$ trend across positions
        \item [GRAPH] cross-correlation NCC + displacement trend across positions
        \item [TEXT] where along the path degradation becomes significant
    \end{itemize}
    \item beam quality
    \begin{itemize}
        \item [GRAPH] summary panel: FWHM, ellipticity, $R^2$, $\omega$ (measured vs. control reference line)
        \item [TEXT] deviation from ideal circular Gaussian, quantified
    \end{itemize}
    \item range angle
    \begin{itemize}
        \item [GRAPH] calibration curve with linear fit (mrad/V)
        \item [TEXT] measured mrad/kV vs. spec, \% agreement
        \item [TEXT] linear range / where nonlinearity appears (if any)
    \end{itemize}
    \item polarization-based angles
    \begin{itemize}
        \item [GRAPH] polarization vs. deflection-axis alignment result
        \item [TEXT] measured axis offset, if any
    \end{itemize}
\end{itemize}

\subsubsection{conclusion}
\begin{itemize}
    \item [TEXT] EOD beam quality: acceptable/degraded range, key limiting factor
    \item [TEXT] deflection linearity and agreement with manufacturer spec
    \item [TEXT] pointing stability / drift implications for MINFLUX acquisition
    \item [TEXT] recommended operating range/voltage for MINFLUX use
    \item [TEXT] open issues or measurements needing follow-up (e.g., polarization/axis alignment)
\end{itemize}
\section{Donut}
\subsection{physical concept}
\begin{itemize}
    \item [TEXT] Laguerre-Gaussian mode, $p=0$, $l=2$ — donut/vortex intensity profile
    \item [TEXT] how EOD/phase element generates the LG mode from the Gaussian input
\end{itemize}

\subsection{characterization}

\subsubsection{intro why}
\begin{itemize}
    \item [TEXT] why the central null matters for MINFLUX (localization precision depends on null depth/sharpness)
    \item [TEXT] limitation cases: imperfect null (extinction ratio), ring asymmetry, deflection-induced null degradation
\end{itemize}

\subsubsection{setup}
\begin{itemize}
    \item donut profile characterization (single frame, fixed position)
    \begin{itemize}
        \item [EXPERIMENT] single donut frame vs. ideal LG model and vs. fitted Gaussian
        \item [GRAPH] X/Y profile: measured vs. LG ideal vs. Gaussian fit
        \item [GRAPH] radial profile (azimuthal average): measured vs. LG ideal vs. Gaussian
        \item [GRAPH] azimuthal profile (ring intensity vs. angle) — ring symmetry
        \item [GRAPH] zoomed 2D heatmap of donut with $r_x$, $r_y$, ER annotated
    \end{itemize}
    \item beam quality (donut)
    \begin{itemize}
        \item extinction ratio (null depth)
        \item ring radius ($r_x$, $r_y$)
        \item beam size: $r_{peak}$, diameter, inner/outer $1/e^2$ radii
        \item LG fit goodness-of-fit ($R^2$)
        \item ring symmetry: coefficient of variation, Fourier amplitudes $a_1$ (tilt), $a_2$ (elliptical), $a_4$ (4-fold), ring ellipticity
    \end{itemize}
    \item range angle (donut deflection sweep)
    \begin{itemize}
        \item [EXPERIMENT] X-axis and Y-axis voltage sweeps, control (no-deflection) reference frame
        \item [GRAPH] extinction ratio vs. voltage (X and Y axes)
        \item [GRAPH] normalized cross-correlation (NCC) vs. voltage (X and Y axes)
        \item [GRAPH] beam center displacement ($\Delta x$, $\Delta y$, px and µm) vs. voltage
        \item [TABLE] summary CSV: max displacement, min/max ER, min/max NCC per axis (with corresponding voltages)
    \end{itemize}
\end{itemize}

\subsubsection{methods}
\begin{itemize}
    \item background subtraction (bottom 5\% pixels)
    \item center detection: centroid of bright ring pixels (30\% peak threshold) — dark null coincides with ring centroid
    \item ring radius: peak location on cross-profile through center, averaged left/right ($r_x$) and top/bottom ($r_y$)

    \item donut profile fitting
    \begin{itemize}
        \item LG ideal model $I(r) = C\,(r/\omega)^4 e^{-2r^2/\omega^2} + bg$, fit by Nelder--Mead least-squares
        \item comparison Gaussian fit (TEM00) centered at same detected null, for reference
        \item $R^2$ computed for both LG and Gaussian fits
    \end{itemize}

    \item extinction ratio (ER)
    \begin{itemize}
        \item ER $= (\text{null min intensity} - bg) / (\text{ring peak intensity} - bg)$
        \item ring peak: mean of top 0.5\% of pixels (robust to hot pixels)
        \item null region: small window at detected center
    \end{itemize}

    \item beam size from radial profile
    \begin{itemize}
        \item $r_{peak}$: radius of peak azimuthally-averaged intensity
        \item inner/outer $1/e^2$ radii: interpolated crossing points on either side of the ring peak
    \end{itemize}

    \item ring symmetry (azimuthal analysis)
    \begin{itemize}
        \item intensity sampled at fixed radius over 360°, averaged over a small radial window
        \item coefficient of variation (uniformity of ring intensity)
        \item Fourier decomposition of azimuthal profile: $a_1$ (dipole/off-center), $a_2$ (elliptical distortion), $a_4$ (4-fold distortion)
        \item ring ellipticity: $\max(r_x,r_y)/\min(r_x,r_y)$
    \end{itemize}

    \item deflection sweep (range angle)
    \begin{itemize}
        \item voltage sweep per axis (X, Y), fixed ROI crop around donut center for each step
        \item center detected independently per frame; displacement relative to control center
        \item ER recomputed per step (same null/ring-peak method, on cropped ROI)
        \item NCC computed against the control ROI via 2D normalized cross-correlation
    \end{itemize}
\end{itemize}

\subsubsection{results}
\begin{itemize}
    \item [GRAPH] measured donut vs. LG ideal vs. Gaussian (profiles + 2D heatmap)
    \item [TEXT] extinction ratio, $R^2$, beam size values at reference (undeflected) position
    \item [GRAPH] ring symmetry metrics (CV, $a_1$, $a_2$, $a_4$, ellipticity) at reference position
    \item [GRAPH] ER and NCC vs. voltage, X and Y axes (deflection-induced degradation)
    \item [GRAPH] displacement vs. voltage, X and Y axes
    \item [TEXT] worst-case ER/NCC and corresponding voltage (from summary CSV)
\end{itemize}

\subsubsection{conclusion}
\begin{itemize}
    \item [TEXT] donut null quality (ER) at rest vs. under deflection
    \item [TEXT] ring symmetry — presence/absence of tilt, ellipticity, or 4-fold distortion
    \item [TEXT] deflection range over which donut quality remains acceptable for MINFLUX
    \item [TEXT] recommended operating voltage range / limitations found
\end{itemize}

\section{Galvometer}
\subsection{physical concept}
\subsection{characterization}
\subsubsection{intro why }
\subsubsection{setup}
\subsubsection{methods}
\subsubsection{results}
\subsubsection{conclusion}
\printbibliography
\section{apendix}




\end{document}
